\section{Centered harmonic moments and the exact Bernoulli transform}
\label{sec:tail-foundation}

We restate the precise inherited moment formulas needed for the new
determinant. This also makes the continuation independent of a
choice of numerical regularization. Fix an even $q=2b\ge4$ and put
\[
 A_s(n)=\zeta(s)-H_n^{(s)}=\zeta(s,n+1),\qquad x_n=n+\frac12,
\]
\begin{equation}\label{eq:tail-Dirichlet}
 D_{2,q}(s)=\sum_{n\ge0} A_2(n)A_q(n)x_n^{-s},
 \qquad \re s>1-q.
\end{equation}
The displayed half-plane is the domain of absolute convergence.
For later $s$, the notation means its unique meromorphic continuation.
We write $D_r=D_{2,q}(-2r)$ for $r\ge0$.

\subsection{Regularity at every required moment}

The centered Hurwitz expansion, obtained from Euler--Maclaurin, is
\begin{equation}\label{eq:centered-hurwitz-foundation}
 A_s(n)\sim x_n^{1-s}\left\{
 \frac1{s-1}+\sum_{h\ge1}
    \frac{B_{2h}(1/2)(s)_{2h-1}}{(2h)!}x_n^{-2h}\right\}.
\end{equation}
Here $s$ in $A_s$ is a fixed positive even integer, and
$(s)_k$ is the rising factorial. The vanishing of all odd
$B_{2h+1}(1/2)$ explains the parity of this expansion.
Consequently $A_2(n)A_q(n)$ has only powers $x_n^{-q-2h}$.
Subtracting finitely many of them from \eqref{eq:tail-Dirichlet}
continues the function to any prescribed left half-plane. Its only
possible poles occur at $s=1-q-2h$, which are negative odd integers.
Every $D_r$ is therefore regular.

This argument also proves a useful operational rule. If $P$ is an
even polynomial and $R$ is the nonnegative-power part of the
asymptotic series for $P(x_n)A_2(n)A_q(n)$, then
\begin{equation}\label{eq:ordinary-fp-foundation}
 \sum_{n\ge0}\{P(x_n)A_2(n)A_q(n)-R(x_n)\}
 =\sum_r [X^{2r}]P(X)\,D_r.
\end{equation}
The left series is absolutely convergent. Each subtracted monomial
has continued sum $\zeta(-2h,1/2)=0$; this follows either from
$\zeta(s,1/2)=(2^s-1)\zeta(s)$ or from the odd Bernoulli polynomial
at $1/2$. Thus the analytic continuation and the explicitly subtracted
ordinary sum have exactly the same constant.

\subsection{The fixed formal coordinate system}

Write decreasing-index multiple zeta values as
\[
 \zeta(a,c)=\sum_{n>m\ge1}\frac1{n^a m^c},\qquad a\ge2,\ c\ge1.
\]
Define the initial blocks by
\begin{align}
 C_0={}&\zeta(2,q-1)+\zeta(q,1)+\zeta(q+1),
 \label{eq:C0-foundation}\\
 C_j={}&\zeta(2,q-2j-1)+\frac12\zeta(q+1-2j)
   +\frac1{2j}\sum_{h=1}^{j}\binom{2j}{2h}
                  B_{2j-2h}\zeta(q-2h),
 \quad 1\le j<b.\label{eq:Cinitial-foundation}
\end{align}
For $j\ge b$, define $C_j=\sum_{k=0}^{b-1}a_{jk}Z_k$, where
\begin{equation}\label{eq:coordinates-foundation}
 Z_0=\zeta(0)=-\frac12,\qquad Z_k=\zeta(2k)\quad(1\le k<b),
\end{equation}
and, with $d=2j+1$, $e=d-q+2k-1$,
\begin{equation}\label{eq:entries-foundation}
 a_{jk}=B_e\left\{
 \frac{\binom{d-q+1}{2-2k}}{d-q+1}
 +\frac{\binom{d-1}{q-2k}}{d-1}\right\}.
\end{equation}
A binomial summand in this formula is omitted unless its lower
index is at least $2$ and no greater than its upper index.
In particular the first summand occurs only when $k=0$.
This is part of the formula, not a convention that may be replaced
by $\binom u0=1$.

\begin{proposition}[Inherited two-tail transform]
\label{prop:tail-transform}
For every $r,j\ge0$,
\begin{align}
 D_r&=\sum_{j=0}^r
       \frac{\binom{2r}{2j}}{2j+1}B_{2r-2j}(1/2)C_j,
       \label{eq:forward-foundation}\\
 C_j&=\sum_{r=0}^j
       \binom{2j+1}{2r}2^{-2(j-r)}D_r.
       \label{eq:inverse-foundation}
\end{align}
Hence all moments lie in the rational span of the $q$ specified
formal labels
\[
 C_0,\ldots,C_{b-1},\qquad Z_0,\ldots,Z_{b-1}.
\]
\end{proposition}

These are the $p=2$ cases of the already proved two-tail formulas
reproduced in \cite[Section 7]{CanonicalFifth}. We give their finite
mechanism to fix the cutoffs and constants. Put
\[
 S_r(N)=\sum_{n=0}^{N-1}(n+1/2)^{2r}
       =\sum_{j=0}^r\frac{\binom{2r}{2j}}{2j+1}
                        B_{2r-2j}(1/2)N^{2j+1}.
\]
This is an odd polynomial with zero constant term. The exact identity
$A_s(k-1)=A_s(k)+k^{-s}$ and finite summation by parts express $D_r$
as the same linear combination of constant terms of
\begin{align}\label{eq:finite-E-foundation}
 E_d(N)={}&N^dA_2(N)A_q(N)
 +\sum_{k=1}^N k^{d-q}A_2(k)
 +\sum_{k=1}^N k^{d-2}A_q(k)
 +\sum_{k=1}^N k^{d-q-2}.
\end{align}
The sign of the first term and the inclusive tail cutoff both
matter. When $d<q$, the first ordered sum tends to
$\zeta(2,q-d)$. For a positive integer exponent $m$, the exact reduction
\begin{equation}\label{eq:finite-Faulhaber-foundation}
 \sum_{k=1}^N k^m A_s(k)
 =A_s(N)F_m(N)+\sum_{\ell=1}^N\frac{F_m(\ell-1)}{\ell^s},
 \quad F_m(N)=\frac{B_{m+1}(N+1)-B_{m+1}}{m+1},
\end{equation}
reduces all remaining sums to finite Bernoulli polynomials and
ordinary power sums. Their constant terms are
\eqref{eq:C0-foundation}--\eqref{eq:entries-foundation}; at the
transition $d=q+1$ the two $-\gamma/2$ contributions cancel the
$+\gamma$ from the final harmonic sum in
\eqref{eq:finite-E-foundation}. The remaining boundary products give
the $Z_0$ column. This is the inherited finite coefficient calculation,
not an assumption of arithmetic independence. Finally the standard
Bernoulli generating identity inverts the odd-polynomial transform,
giving \eqref{eq:inverse-foundation}.

The existing report proves that a longer specified moment prefix
has full formal rank. The next section determines the exact first-$q$
minor for every $p=2$, without any numerical nonvanishing assumption.
